\section{Discussion}

A reading panel of nine readers built from one VLM did not provide nine independent reads. Across the original-image specialty panels, error correlation ranged from 0.685 to 0.836 for GPT-5 mini, Gemini 3.7 Flash, and Claude Haiku 4.5. The effective number of independent readers ranged from 1.17 to 1.39. Qwen 3.8 Flash had lower error correlation, 0.420--0.651, and an effective number of independent readers of 1.45--2.06. The reading panels, discussion, and panel arithmetic in Fig. 2 explain the limited return from adding readers. Across every nine-reader panel on yes/no items, observed majority accuracy closely matched the beta-binomial prediction from mean reader accuracy and error correlation. The observed-minus-predicted difference had a mean of -0.12 percentage points and an SD of 0.84 percentage points. Majority voting changed accuracy by -2.1 to +3.5 percentage points relative to the panel's own mean reader. Qwen's larger gain over the single reader came mainly from higher reader accuracy. On VQA-RAD, its mean panel reader exceeded the single reader by +9.7 percentage points (95\% CI, +4.6 to +15.0). Aggregation added +3.5 percentage points (95\% CI, +1.0 to +6.1). Reader accuracy and aggregation are distinct sources of improvement.

Discussion converged on agreement rather than consistently correcting errors. Final error correlation reached 0.843--1.000, and unanimity reached 85.2\%--100.0\%, yet 8.9\%--70.0\% of unanimous panels were wrong. The discussion and conformity analyses in Fig. 9 show how correct answers were lost during this convergence. Correct answers that differed from the gray-image answer were abandoned more often than correct answers that matched it, especially on MedXpertQA-MM. These image-dependent correct answers were also more often held by a minority on that set. Discussion therefore exposed the answers most dependent on image evidence to pressure from a majority sharing the question-text prior. Minority correct answers were overturned in most model--data set combinations, whereas majority correct answers were usually retained. Agreement after discussion consequently reflects both initial evidence and deference to other readers. In mammography, radiologists followed incorrect AI BI-RADS suggestions, demonstrating deference even when the suggested assessment was wrong\cite{dratsch2023automation}. In the reading panels studied here, agreement among readers who share a prior carries that prior forward. Unanimity after discussion is therefore insufficient evidence of a correct reading.

The question-text prior was a common cause of shared errors. The graded-evidence experiment in Fig. 3 showed that removing the image lowered accuracy toward chance on VQA-RAD and SLAKE. On MedXpertQA-MM, it changed single-reader accuracy little for GPT-5 mini, Claude Haiku 4.5, and Qwen 3.8 Flash. Without the image, error correlation rose for Gemini, Claude Haiku, and Qwen on VQA-RAD and SLAKE: readers converged on answers supplied by the remaining text. The reader-level decomposition in Fig. 4 connected these answers to errors with real images. Among GPT-5 mini, Gemini, and Claude Haiku, 33.9\%--72.4\% of consensus errors repeated the gray-image majority; Qwen's range was 15.0\%--77.5\%. Many reader answers were wrong both with and without the image, particularly on MedXpertQA-MM. Language priors in general visual question answering, poor performance under medical probing, and limited sensitivity to image substitution provide related evidence\cite{goyal2017vqa,yan2025worse,zafar2026see}. Clinical information also changes how diagnostic tests are read\cite{loy2004clinical}. Here, the question text acts like clinical information that steers the reading. When it favors an incorrect answer, readers sharing the model also share a route to that error.

Readers used the image unevenly, even when they answered the same kinds of questions. Region occlusion in Fig. 5 changed GPT-5 mini answers substantially more for the queried region than for a matched control region. Its panel-reader region dependence was +17.1 percentage points (95\% CI, +13.2 to +21.2), compared with +3.2 percentage points (95\% CI, +0.8 to +5.7) for Claude Haiku. Open-weight models ranged from +2.4 to +21.6 percentage points. The image pairs in Fig. 5 tested whether the same question elicited opposite answers when the images required them. GPT-5 mini and Claude Haiku single readers gave different answers on only 42.5\%--62.0\% of pairs; their panel readers did so on 45.3\%--66.0\%. Open-weight models distinguished the paired images on 39.5\%--63.5\% of pairs. The remaining pairs received the same answer despite opposite reference answers. These interventions establish image use more directly than correctness on isolated questions. Accuracy alone overstates image use when a correct answer is supported by the text, consistent with reports of correct medical answers accompanied by flawed image comprehension\cite{jin2024hidden}.

Seeing set the ceiling for judging from reports. In Fig. 6, question-blind reports written by Gemini 3.7 Flash raised accuracy for both complete reasoners, GPT-5 mini and Claude Haiku 4.5, on every set. Each reasoner's own reports did not consistently improve on reading the image directly. Separating seeing from the question therefore helped when the perceiver supplied better findings, rather than simply because an intermediate report was introduced. Readers sharing one report agreed almost completely on VQA-RAD and SLAKE, with error correlation of 0.944--0.986. Fixing the report removed most variation among readers of the same model, locating that variation primarily in seeing. On MedXpertQA-MM, shared-report error correlation was lower, 0.800--0.893, and the reasoner also mattered. Given the same Gemini reports, GPT-5 mini reached 62.2\% accuracy (95\% CI, 57.2\% to 67.2\%), whereas Claude Haiku reached 45.6\% (95\% CI, 40.0\% to 50.9\%). Better findings improved both reasoners, but identical findings did not produce identical accuracy. The report pathway thus separates the contribution of extracted image findings from the reasoner's ability to select an answer using those findings.

Independent second reads require differences in what readers see or in the models that process the evidence. In Fig. 7, the mean contrasts over models lowered error correlation by small amounts for roles (-0.059 to -0.044), renderings (-0.128 to -0.069), and private reference cases (-0.166 to -0.146). A different perceiver lowered it more (-0.332 to -0.260), and a different model lowered it most (-0.583 to -0.487). These are changes in error correlation relative to the matched shared-component conditions. Lower error correlation alone did not guarantee higher panel accuracy. In Fig. 8, Claude Haiku's SLAKE panels with shared versus private text reference cases had error correlations of 0.941 (95\% CI, 0.907 to 0.970) and 0.720 (95\% CI, 0.657 to 0.779). Accuracy remained 78.0\% in both conditions (95\% CIs, 73.2\% to 83.2\% and 73.2\% to 82.8\%, respectively). For radiology practice, a panel of one model's readers should be counted as one read. An AI second read requires a different model or a different perceiver to supply additional independence. Human evidence likewise supports pooling independent judgments of similar accuracy\cite{kurvers2016pooling}, and independent double reading improves mammographic cancer detection\cite{taylor2008double}. The relevant property is complementary errors at adequate reader accuracy.

The study measured answers to closed questions rather than clinical reporting. Items were yes/no or multiple-choice questions from public data sets. Each condition used one sampling seed, and the confidence intervals resampled items; they did not measure variability across sampling seeds. API models were evaluated at one time point through one router. Open-weight models had four to eight billion parameters. No human readers participated. Gray-image abstentions counted as wrong, so refusal to answer contributed to measured error. Error correlation on items answered by both readers was included as a sensitivity analysis.
